% Chapter 4 - The McEliece Cryptosystem
\section{The McEliece Cryptosystem}
\label{chap:mceliece}

\subsection{Error-Correcting Codes}
\label{sec:errorCorrectingCodes}
% TODO: siehe Strukturplan.tex, Abschnitt 4.1
% - Motivation, Kommunikationsfehler, Fehlerkorrektur

\subsection{Linear Codes}
\label{sec:linearCodes}
% TODO: siehe Strukturplan.tex, Abschnitt 4.2
% - Codewords, Generator Matrix, Parity Check Matrix

\subsection{Goppa Codes}
\label{sec:goppaCodes}
% TODO: siehe Strukturplan.tex, Abschnitt 4.3
% - Grundidee, Eigenschaften, Fehlerkorrekturfaehigkeit

\subsection{Structure of the McEliece Cryptosystem}
\label{sec:mcelieceStructure}
% TODO: siehe Strukturplan.tex, Abschnitt 4.4
% - Private Key, Public Key, Message, Error Vector

\subsection{Key Generation}
\label{sec:keyGeneration}
% TODO: siehe Strukturplan.tex, Abschnitt 4.5
% - Generator Matrix, Scrambling Matrix, Permutation Matrix, Public Key Construction

\subsection{Encryption}
\label{sec:mcelieceEncryption}
% TODO: siehe Strukturplan.tex, Abschnitt 4.6
% - Encoding, Error Addition

\subsection{Decryption}
\label{sec:mcelieceDecryption}
% TODO: siehe Strukturplan.tex, Abschnitt 4.7
% - Decoding, Error Correction

\subsection{Security Assumption}
\label{sec:securityAssumption}
% TODO: siehe Strukturplan.tex, Abschnitt 4.8
% - General Decoding Problem, Schwierigkeit zufaelliger linearer Codes

\subsection{Educational Abstraction}
\label{sec:educationalAbstraction}
% TODO: siehe Strukturplan.tex, Abschnitt 4.9
